%============================================================
\chapter{無人化：運営者を除いた極限}
\label{part:unmanned}
%============================================================

\section{問い}

第\ref{ch:solo}章は一人事業を扱った。本章はその極限、すなわち
運営者が存在しない事業が成立するかを問う。
立ち上げに人手が要ることは自明であるため、
定常運転に限って無人化が可能かを検討し、次いで立ち上げ費用そのものを
どこまで削減できるかを扱う。


\begin{center}
\begin{tabularx}{\textwidth}{@{}llX@{}}
\toprule
\multicolumn{3}{@{}l}{\textbf{本章で扱う理論量}}\\
\midrule
$\Omega,\ \hat\Omega$ & 第\ref{sec:info}章 & 状態空間と、設計時に想定した部分集合 \\
$\lambda_{\rm novel}$ & 本章 & $\Omega\setminus\hat\Omega$ の到来率 \\
$\mathcal{G}$ & 第\ref{sec:info}章 & 自動化はこれより強い条件を要求する \\
$\kappa_C>0$ & 式\eqref{eq:kappa} & 無料期間として自ら発行する信用 \\
\bottomrule
\end{tabularx}
\captionof{table}{本章で扱う理論量の一覧。}
\label{tab:unmanned-quantities}
\end{center}

\section{無人の定式化}

\subsection{自動化は可測性より強い条件である}

第\ref{sec:info}章では、契約が実装可能である条件を $\mathcal{G}$ 可測性とした。
自動化はこれより強い。

\begin{definition}[自動化]
応答を状態空間上の可測関数 $a:\Omega\to A$ として\emph{事前に}書き下すこと。
\end{definition}

命題\ref{prop:implementable}は $\pi$ が $\mathcal{G}$ 可測であることを要求した。
自動化はこれに加えて、応答が\emph{事前に}書き下されていることを要求する。
契約は事後に検証できれば足りるが、自動化はそうではない。
アルゴリズムとは、事前に指定された可測関数にほかならない。
すなわち
\begin{equation}
\text{実装可能}\;:\;\pi\ \text{が}\ \mathcal{G}\ \text{可測}
\qquad\subsetneq\qquad
\text{自動化可能}\;:\;a\ \text{が}\ \hat\Omega\ \text{上で事前指定済}
\end{equation}
という包含関係が成立する。

問題は $\Omega$ が閉じていないことである。
設計時に想定した状態空間を $\hat\Omega\subset\Omega$ とすると、
$\omega\in\Omega\setminus\hat\Omega$ が到来したとき $a$ は定義されていない。
法改正、脆弱性の発見、決済事業者の規約変更、想定外の使用がこれにあたる。

\begin{proposition}[無人化は期間の関数である]
\label{prop:unmanned-period}
未想定状態の到来率を $\lambda_{\rm novel}$ とすると、
無人で運転できる期間はおよそ $1/\lambda_{\rm novel}$ である。
有限期間の完全無人は成立するが、無期限には成立しない。
\end{proposition}

\subsection{必要なのは労働ではなく可用性}

$\Omega\setminus\hat\Omega$ への応答は定常的な労働ではない。
通常は何も行わず、到来したときにのみ動く。
この構造は第\ref{part:catalog}部の族2-5（オプション・保証）と同型である。

\begin{equation}
\text{無人事業の運営者は }\Omega\text{ に対するショートオプションを保有する}
\end{equation}

平常時の履行 $\delta$ はゼロにできるが、\textbf{可用性はゼロにできない}。
労働量は限りなく減らせるが、待機は残る。
「無人」と呼ばれるものは、フローとしての労働が無人なのであって、
オプションとしての可用性は有人である。

これは第\ref{part:solo}章の容量制約とは別種の拘束である。
式\eqref{eq:solo-capacity}が履行の\emph{総量}に上限を課すのに対し、
可用性は\textbf{連続して応答できない区間の長さ}に上限を課す。

\begin{definition}[応答期限と可用性制約]
\label{def:availability}
$\omega\in\Omega\setminus\hat\Omega$ の到来から応答までに許される時間を
\emph{応答期限}と呼び $T_{\rm resp}$ と書く。
運営者が連続して応答できない区間の長さを $u$ とすると、可用性制約は
\begin{equation}
\sup u \;\leq\; T_{\rm resp}
\label{eq:availability}
\end{equation}
と書ける。
\end{definition}

式\eqref{eq:solo-capacity}と式\eqref{eq:availability}は次元がともに時間であるが、
拘束する対象が異なる。前者は積分に、後者は上限に効く。
\textbf{したがって労働量をいくら減らしても式\eqref{eq:availability}は緩まない。}
一人事業の実務的負担がここにあるという説明が可能である。

\begin{remark}[$T_{\rm resp}$ の水準は成分ごとに異なる]
\label{rem:tresp-components}
$T_{\rm resp}$ は $\Omega\setminus\hat\Omega$ の内容に依存する。
第\ref{ch:platformfee}章で $\lambda_{\rm novel}$ を四成分に分解するが、
脆弱性の公表と法改正では猶予の桁が異なる。
式\eqref{eq:availability}を実際に評価するには成分ごとの $T_{\rm resp}$ を要するが、
本稿はこれを測定していない。
\end{remark}

\begin{remark}[ナイトの区別との対応]
\label{rem:knight}
$\hat\Omega$ 上の変動は測度が定義されているためリスクであり、
自動化も保険（族5-4）も可能である。
$\Omega\setminus\hat\Omega$ は不確実性であり、どちらもできない。
\end{remark}

\begin{remark}[無人化の程度は定義できない]
\label{rem:automation-rate-illposed}
「どの程度無人か」を比として表すことはできない。

履行の側で測る経路は成立しない。
$\delta$ は評価写像 $v$ により貨幣単位を持つが、
$v$ が測るのは価値であって手間ではない。
自動化された履行と人手による履行の価値の比は、無人化の程度を表さない。

応答の側で測る経路も成立しない。
自動化とは $a$ が $\hat\Omega$ 上で事前指定されていることであるから、
程度は $\hat\Omega$ と $\Omega$ の大きさの比になる。
しかし注\ref{rem:knight}のとおり $\Omega\setminus\hat\Omega$ には測度が定義されない。

\textbf{したがって障害は単位の不在ではなく測度の不在であり、
定義を精密化しても解消しない。}
無人化は程度ではなく期間の関数として述べるほかない（命題\ref{prop:unmanned-period}）。
\end{remark}

\begin{remark}[無人化と利潤の関係]
$\Phi$ と $f$ が完全に事前指定できるならば、それは完全に複製可能である。
したがって競争により $\phi_{\rm prod}\to 0$ に向かう圧力が働く。
完全無人に近づくほど利潤の源泉が薄くなる。

利潤とは自動化できない残余への報酬である、という読み方が可能である。
ネットワーク効果やブランドが例外を作るため断定はできないが、
機構としては働いていると考えられる。
\end{remark}

\subsection{無料期間の位置づけ}

無料試用期間は、$\delta$ を提供して $\pi=0$ とする状態、
すなわち事業者が顧客に与信する $\kappa>0$ である。
これが成立するのは\textbf{限界費用がゼロであるため与信コストがほぼゼロ}だからである。
物理財の無料試用は在庫を失うが、ソフトウェアの無料期間はインフラ費のみで済む。

\begin{equation}
\text{無料トライアル} = \text{限界費用ゼロの事業のみが安価に発行できる信用}
\end{equation}

第\ref{sec:bootstrap}節で信用のブートストラップ問題の解を三つ挙げたが、
これは第四の解である。すなわち\textbf{自ら与信して信用を代替する}。
$T_{\rm credit}$ を $\kappa>0$ のコストで購入していることになる。

\section{ブートストラップの分解}

\subsection{各項と外部化先}

式\eqref{eq:E-required-2}の各項について、誰が吸収しうるかを整理する。
制度層への登録に要する期間を $T_{\rm inst}$ として追加する。

\begin{center}
\begin{tabularx}{\textwidth}{@{}lXXl@{}}
\toprule
項 & 内容 & 外部化先 & 対価 \\
\midrule
$T_{\rm dev}$ & 履行系の構築 & 既製基盤、決済SDK、ホスティング & 利用料 \\
$T_{\rm inst}$ & 制度層への登録 & 決済代行の加盟店審査、MoR & 手数料 \\
$T_{\rm credit}$ & 信用の蓄積 & プラットフォームの返金保証、無料期間で自弁 & 手数料、$\kappa>0$ \\
$T_{\rm acq}$ & 顧客の獲得 & ストア、マーケットプレイスの流通 & 販売手数料 \\
$\hat\Omega$ の列挙 & 想定状態の設計 & \textbf{不可} & --- \\
$\Phi$ の選択 & 契約形態の決定 & \textbf{不可} & --- \\
\bottomrule
\end{tabularx}
\captionof{table}{ブートストラップ費用の各項と外部化先、対価。}
\label{tab:bootstrap-outsourcing}
\end{center}

\begin{proposition}[削減不能な残余]
外部化できないのは $\Phi$ の選択と $\hat\Omega$ の列挙の二つのみである。
\end{proposition}

外部化には一貫した代償がある。
\textbf{$\phi_{\rm barg}$ を恒久的に手放すこと}である。
一回限りの立ち上げ費用を、手数料という永続的な費用に変換している。

\subsection{族ごとの主成分}

第\ref{part:solo}章で生存した類型について、立ち上げ費用の主成分が異なる。

\begin{center}
\begin{tabular}{@{}llllll@{}}
\toprule
類型 & $T_{\rm dev}$ & $T_{\rm credit}$ & $T_{\rm acq}$ & 主成分 \\
\midrule
1-1 デジタル商品の売切り & 高 & 低 & 低 & $T_{\rm dev}$ \\
2-1 定額アクセス & 高 & 中 & 中 & 分散 \\
6-4 機会課金 & 中 & 低 & 高 & $T_{\rm acq}$（両面市場） \\
7-4 寄付・支援 & 低 & 高 & 中 & $T_{\rm credit}$ \\
5-5 IPライセンス & 中 & 高 & 中 & $T_{\rm credit}$ \\
\bottomrule
\end{tabular}
\captionof{table}{生存類型ごとの立ち上げ費用の主成分。}
\label{tab:startup-cost-by-type}
\end{center}

1-1 と 7-4 の主成分が互いに補完的である点に注目すると、
族間の遷移として立ち上げ費用を分割する設計が導かれる。

\begin{center}
\begin{tabularx}{\textwidth}{@{}lX@{}}
\toprule
段階 & 内容 \\
\midrule
公開 & $T_{\rm credit}$ を積む。収益ゼロ。副産物として $T_{\rm dev}$ が進行する \\
7-4 支援 & 積んだ信用のみで最初の $\pi$ が立つ。$T_{\rm dev}$ をほぼ要求しない \\
1-1 売切り & $T_{\rm credit}$ が済んでいるため立ち上げが軽い。$\kappa\approx 0$ \\
2-1 定額 & $\kappa<0$ となり、第\ref{sec:growth}章の意味で成長が現金を生む \\
\bottomrule
\end{tabularx}
\captionof{table}{族間遷移による立ち上げ費用の分割設計。}
\label{tab:phase-transition-design}
\end{center}

第\ref{part:solo}章の漏れ出し項 $\theta$（式\eqref{eq:leakage}）が、
ここでは公開から $T_{\rm dev}$ への寄与として逆向きに働く。
\textbf{立ち上げ費用は一括で支払う必要がなく、族を渡り歩くことで分割できる。}

\section{本章の限界}
\label{sec:unmanned-limits}

\begin{enumerate}[label=(\arabic*)]
\item \textbf{無人化の程度を表す量を持たない。}
注\ref{rem:automation-rate-illposed}のとおり、障害は単位の不在ではなく測度の不在であり、
定義の精密化では解消しない。本章が述べるのは期間についてのみである。

\item \textbf{$T_{\rm resp}$ の水準を与えていない。}
式\eqref{eq:availability}は書けるが、その右辺は
$\Omega\setminus\hat\Omega$ の成分ごとに異なり、本章は値を特定しない
（注\ref{rem:tresp-components}）。

\item \textbf{無人化と利潤の逆相関は機構の指摘にとどまる。}
反例（ネットワーク効果、ブランド）の扱いを含めた検証を行っていない。

\item \textbf{立ち上げ費用の外部化価格を持たない。}
表\ref{tab:bootstrap-outsourcing}は外部化先を列挙するが対価を数値で与えていない。
第\ref{ch:platformfee}章でこれを測定する。
\end{enumerate}

